\documentclass[11pt,a4paper]{article}

\usepackage[margin=1in]{geometry}
\usepackage{amsmath,amssymb,amsthm}
\usepackage{booktabs}
\usepackage{tikz}
\usepackage{listings}
\usepackage{xcolor}
\usepackage{hyperref}
\usepackage{graphicx}
\usepackage{microtype}
\usepackage{cite}

\hypersetup{
  colorlinks=true,
  linkcolor=blue!50!black,
  citecolor=blue!50!black,
  urlcolor=blue!70!black
}

\newtheorem{theorem}{Theorem}
\newtheorem{lemma}{Lemma}
\newtheorem{definition}{Definition}
\newtheorem{proposition}{Proposition}

\lstset{
  basicstyle=\ttfamily\small,
  breaklines=true,
  frame=single,
  columns=fullflexible,
  keepspaces=true,
}

\title{ZAP-FIX: A Zero-Copy, Post-Quantum Transport for the Financial Information eXchange Protocol}

\author{ZAP Protocol Authors\\\texttt{zap-proto.io}}

\date{\today}

\begin{document}

\maketitle

\begin{abstract}
The Financial Information eXchange (FIX) protocol is the dominant
session-layer protocol for institutional equities, futures, and
derivatives order flow. Its standard transport binding (raw TCP, or
TCP with TLS~1.3) was designed in an era when (a) bandwidth was the
binding constraint, (b) nanosecond-scale latency was a curiosity rather
than a regulatory expectation, and (c) cryptographically relevant
quantum adversaries were a long-horizon concern. None of these
conditions hold for institutional venues today. We present
\emph{ZAP-FIX}, a transport binding for FIX 4.4 and FIX 5.0 SP2 that
preserves the FIX session layer in full---SequenceNumber discipline,
Logon/Logout, GapFill, ResendRequest, Heartbeat---while replacing the
underlying TCP/TLS stream with the ZAP zero-copy, mutually
authenticated, post-quantum frame protocol. We give a formal latency
model, a compatibility analysis, a resilience model, and a compliance
mapping to MiFID~II RTS~25 timestamping. Our analysis shows that
ZAP-FIX delivers steady-state per-message latency strictly less than
or equal to FIX-over-TLS, with the gap dominated by the elimination
of SOH~(\texttt{0x01}) text scanning in favour of zero-copy schema
fields, and that session-resumption latency is bounded by a single
ZAP handshake even after pod migration or network partition---a regime
in which TCP+TLS today requires connection re-establishment plus
session-state reconstruction.
\end{abstract}

\section{Introduction}
\label{sec:introduction}

The Financial Information eXchange (FIX) protocol, in versions 4.2,
4.4, and 5.0~SP2, is the de~facto session protocol for institutional
order flow across cash equities, listed futures, options, and a large
fraction of OTC derivatives~\cite{fix44,fix50sp2}. Conservative public
estimates put the share of institutional equity order messages carried
over a FIX session at well above ninety percent. The protocol layers
cleanly over a reliable byte-stream transport: implementations
overwhelmingly choose TCP, with TLS~1.3 layered in for venues subject
to confidentiality or integrity regulation~\cite{rfc8446,rfc9293}.

Two pressures are now reshaping the transport choice. First,
deterministic tail latency is an explicit best-execution
obligation under MiFID~II RTS~25, which requires
microsecond-granular timestamps for many asset
classes~\cite{rts25}. Co-located institutional clients today operate
inside a round-trip latency budget of approximately
$50$--$200~\mu\mathrm{s}$ from order entry to acknowledgement,
with tail latencies dominated not by steady-state per-message cost
but by TLS handshake re-establishment and FIX session-state
reconstruction following a transient
disconnection~\cite{hasbrouck2013,kissell2013}. Second, the prospect
of a cryptographically relevant quantum adversary has motivated
regulators and venue operators to require, or at least to permit,
post-quantum confidentiality for stored and in-flight order flow,
under a ``harvest-now, decrypt-later'' threat
model~\cite{nistpqc,nistsp80056c}.

\paragraph{Contribution.}
We define \emph{ZAP-FIX}: a transport binding that preserves the FIX
session protocol as specified by FIX Protocol Limited and
substitutes only the byte-stream layer. Concretely:

\begin{itemize}
  \item Each FIX administrative or application message becomes
        exactly one ZAP frame. There is no in-flight chunking and
        no aggregation of multiple FIX messages into a single
        frame, preserving message-boundary semantics.
  \item The SOH-delimited \texttt{tag=value} body is wrapped in a
        Cap'n Proto envelope; a small, fixed set of routing fields
        (\texttt{SenderCompID}, \texttt{TargetCompID},
        \texttt{MsgSeqNum}, \texttt{MsgType}) is lifted into typed
        envelope slots so that the gateway can dispatch without
        scanning the FIX body.
  \item Session establishment uses the ZAP handshake, which provides
        mutual authentication and X-Wing hybrid post-quantum key
        agreement~\cite{xwing,kyber} in a single round trip.
  \item SequenceNumber discipline is preserved exactly: ZAP frame
        ordering on a session is constrained to match FIX
        \texttt{MsgSeqNum} monotonicity, so \texttt{ResendRequest}
        and \texttt{GapFill} semantics remain intact.
  \item An optional \emph{ZAP-FIX BinaryEncoding} mode bit-packs
        frequent typed fields (\texttt{Symbol}, \texttt{Price},
        \texttt{OrderQty}) directly into the envelope, reducing
        on-wire bytes by 30--50\% relative to SOH text and
        eliminating the corresponding format/parse cost. This mode
        is conceptually adjacent to FIX SBE~\cite{fixsbe} and FIX
        FAST~\cite{fixfast}, but it inherits the ZAP framing and
        cryptographic guarantees rather than layering separately.
\end{itemize}

We also specify a transparent ZAP-FIX gateway, which terminates
legacy TCP/TLS-FIX on the front side and speaks ZAP-FIX on the
back side, so that institutional clients carrying long-lived
connectivity contracts and certifications can adopt ZAP-FIX
incrementally without endpoint migration.

The remainder of the paper is organised as follows.
Section~\ref{sec:fix-recap} recaps the FIX session layer.
Section~\ref{sec:construction} specifies the ZAP-FIX construction.
Section~\ref{sec:latency} gives the latency analysis.
Section~\ref{sec:resilience} describes the resilience model.
Section~\ref{sec:compliance} maps the construction to MiFID~II RTS~25
and FINRA~OATS audit
obligations~\cite{rts25,finra-oats,finra-cat}.
Section~\ref{sec:compat} discusses compatibility with deployed FIX
infrastructure. Section~\ref{sec:related} surveys related work.

\section{FIX Session Layer Recap}
\label{sec:fix-recap}

We adopt the standard FIX session-layer
abstractions~\cite{fix44,fix50sp2}.

A \emph{session} is a bidirectional, sequenced exchange of FIX
messages between two parties identified by \texttt{SenderCompID} and
\texttt{TargetCompID}. Each direction maintains a strictly
increasing \texttt{MsgSeqNum}. The protocol is intentionally
stream-oriented: the underlying transport is required to deliver
bytes in order and to surface gaps, but is permitted to coalesce or
fragment messages at byte granularity. The session layer compensates
with a small set of administrative messages.

\paragraph{Logon (\texttt{A}).}
Initiates a session, negotiates encryption (historically; in modern
deployments \texttt{EncryptMethod}~$=0$ and confidentiality is
provided by the transport), heartbeat interval, and optional
sequence-number reset.

\paragraph{Heartbeat (\texttt{0}) and TestRequest (\texttt{1}).}
Liveness probes; absence of either over an interval triggers session
termination.

\paragraph{ResendRequest (\texttt{2}) and SequenceReset (\texttt{4}).}
The receiver, on detection of a gap (\texttt{MsgSeqNum} greater than
expected), issues \texttt{ResendRequest} for the missing range; the
sender replays application messages and may emit a
\texttt{SequenceReset}--GapFill for administrative messages it
declines to replay.

\paragraph{Logout (\texttt{5}).}
Graceful termination.

The transport contract that the FIX session layer requires is
therefore quite narrow: in-order, reliable, byte-oriented delivery
within a single connection lifetime, with explicit failure (the
connection is torn down) when delivery cannot be guaranteed. Any
transport satisfying this contract is admissible. ZAP, restricted
to a single-session ordered binding, satisfies it; we make this
precise in Section~\ref{sec:construction}.

\section{ZAP-FIX Construction}
\label{sec:construction}

\subsection{Frame Mapping}

Let $M$ be a FIX message with header fields
$h = (\texttt{SenderCompID}, \texttt{TargetCompID},
\texttt{MsgSeqNum}, \texttt{MsgType})$
and body $b$ (the \texttt{tag=value} fields excluding $h$, with the
SOH-delimited textual encoding $\sigma(M)$).

The ZAP-FIX frame for $M$ is the Cap'n Proto envelope
\[
F(M) \;=\; \langle\, \mathrm{hdr}(h),\; \mathrm{body}(M)\, \rangle
\]
where $\mathrm{hdr}(h)$ is a fixed-layout typed struct exposing the
four routing fields, and $\mathrm{body}(M)$ is one of the following
two encodings, negotiated at session \texttt{Logon}.

\begin{itemize}
  \item \textbf{Text mode.} $\mathrm{body}(M) = \sigma(M)$. The
        SOH-delimited canonical FIX body is carried verbatim as a
        single byte slice. Conformance to the FIX specification is
        bit-for-bit on the body; the envelope is metadata only.
  \item \textbf{BinaryEncoding mode.} A small dictionary of typed
        fields (\texttt{Symbol}, \texttt{Price}, \texttt{OrderQty},
        \texttt{Side}, \texttt{OrdType}, \texttt{TimeInForce},
        \texttt{TransactTime}, \texttt{ClOrdID}) is bit-packed in
        the envelope. Remaining fields fall through to a residual
        SOH body. Decoders MUST be capable of reconstructing the
        canonical SOH form from the envelope. This mirrors the
        intent of FIX~SBE~\cite{fixsbe} but does not require a
        separate transport binding.
\end{itemize}

\paragraph{One-message-per-frame invariant.}
For every FIX message $M$, exactly one ZAP frame $F(M)$ is emitted on
the session, and ZAP frames are not coalesced. This is the only
deviation from a faithful byte-stream transport, and it strengthens,
rather than weakens, the FIX session contract: gaps are detectable at
frame granularity rather than at byte granularity.

\subsection{Session Establishment}

ZAP-FIX session establishment is a strict superset of the FIX Logon
exchange:

\begin{enumerate}
  \item \textbf{ZAP handshake.} A single round trip exchanges static
        and ephemeral X25519 public keys, X-Wing
        encapsulations~\cite{xwing,kyber}, and a transcript hash. On
        success, both parties hold a shared traffic secret and a
        peer identity bound to a static long-term key.
  \item \textbf{FIX Logon~(\texttt{A}).} The first FIX frame
        carries Logon. \texttt{EncryptMethod} is fixed to $0$
        (transport provides confidentiality and integrity).
        \texttt{HeartBtInt}, \texttt{ResetSeqNumFlag}, and
        \texttt{NextExpectedMsgSeqNum} retain their FIX semantics.
\end{enumerate}

The handshake binds the FIX \texttt{SenderCompID} to the ZAP
identity. The gateway rejects the Logon if the asserted CompID does
not appear in the entitlement list authorised for the ZAP identity,
defending against a class of CompID-spoofing attacks that pure
TCP/TLS-FIX cannot detect at the transport layer.

\subsection{Sequence-Number Discipline}

The ZAP-FIX session maintains, for each direction, the FIX
invariant
\begin{equation}
  \texttt{MsgSeqNum}(F_{i+1}) = \texttt{MsgSeqNum}(F_i) + 1.
  \label{eq:seqnum}
\end{equation}
Because frames are not coalesced and ZAP delivery is in-order on a
session, the receiver detects a gap by inspecting the envelope
header alone, without parsing the SOH body. This is the basis for
the parse-cost advantage of Section~\ref{sec:latency} and for
faster \texttt{ResendRequest} of Section~\ref{sec:resilience}.

\section{Latency Analysis}
\label{sec:latency}

We decompose the per-message round-trip latency between two co-located
peers into:
\begin{equation}
  T_{\mathrm{rtt}} \;=\; T_{\mathrm{ser}} + T_{\mathrm{tx}}
                       + T_{\mathrm{net}}
                       + T_{\mathrm{rx}} + T_{\mathrm{parse}}
                       + T_{\mathrm{app}} + T_{\mathrm{rev}}
  \label{eq:budget}
\end{equation}
where $T_{\mathrm{ser}}$ is sender-side serialisation,
$T_{\mathrm{tx}}$ kernel transmit, $T_{\mathrm{net}}$ wire and
switch fabric, $T_{\mathrm{rx}}$ kernel receive, $T_{\mathrm{parse}}$
receiver-side parse, $T_{\mathrm{app}}$ application processing
(matching, risk, etc.), and $T_{\mathrm{rev}}$ the reverse-path
contributions.

The transport choice fixes
$T_{\mathrm{ser}}$, $T_{\mathrm{parse}}$, and partially
$T_{\mathrm{net}}$ (via on-wire byte count). $T_{\mathrm{app}}$ is
unaffected.

\subsection{Steady-State Per-Message Cost}

\paragraph{Serialisation.}
For FIX-over-TCP/TLS, $T_{\mathrm{ser}}$ is dominated by formatting
typed fields into ASCII and inserting SOH separators; in optimised
implementations this lands in the $1$--$5~\mu\mathrm{s}$ range per
message. For ZAP-FIX text mode, the cost is the same plus a
constant Cap'n Proto envelope ($< 100~\mathrm{ns}$ on contemporary
hardware). For ZAP-FIX BinaryEncoding mode, formatting of the
hot fields is replaced by bit-packing of pre-typed values
(approximately $50$--$200~\mathrm{ns}$ for the hot field set).

\paragraph{Parse.}
For FIX-over-TCP/TLS, $T_{\mathrm{parse}}$ is dominated by an SOH
scan and \texttt{tag=value} tokenisation; routing fields are
reachable only after the scan reaches them. For ZAP-FIX, routing
fields are zero-copy reads from the envelope; an end-to-end FIX
session that only requires the routing fields plus a handful of
hot fields can avoid touching the SOH body altogether.

\paragraph{On-wire bytes.}
A typical \texttt{NewOrderSingle} in SOH text is in the range
$180$--$280$ bytes. ZAP-FIX BinaryEncoding mode reduces this to
$110$--$170$ bytes for the same field set. Within a single
$1500$-byte MTU, $T_{\mathrm{net}}$ is unaffected; over multi-segment
or constrained fabrics, the byte reduction translates linearly.

We make this precise in Theorem~1 of the companion proof document.

\subsection{Handshake and Resumption}

A TLS~1.3 handshake completes in one round trip in the
\texttt{0-RTT} or \texttt{1-RTT} regime~\cite{rfc8446}. The ZAP
handshake also completes in one round trip and yields the same
forward-secrecy and authentication properties, plus post-quantum
KEM hybridisation~\cite{xwing}. In steady state the handshake cost
is amortised; the regime in which it dominates is connection
re-establishment, which we treat in Section~\ref{sec:resilience}.

\section{Resilience}
\label{sec:resilience}

A FIX session over TCP+TLS that loses its underlying TCP connection
must:
\begin{enumerate}
  \item Re-establish a TCP connection (one round trip plus retransmit
        backoff if loss was the cause).
  \item Re-establish a TLS session (one round trip; or
        \texttt{0-RTT} resumption with replay-window caveats).
  \item Reissue \texttt{Logon}.
  \item Detect the gap by comparing \texttt{MsgSeqNum} expectations.
  \item Issue \texttt{ResendRequest} and replay the missing range.
\end{enumerate}
Empirically the dominant cost in steps 4--5 is the byte-stream scan
of the receive buffer to locate FIX message boundaries, since the
TCP byte-stream surface does not expose them.

ZAP-FIX collapses this:
\begin{enumerate}
  \item One ZAP handshake re-establishes the secure channel.
  \item \texttt{Logon} reasserts the session.
  \item Gap detection is a single integer comparison on the next
        envelope header.
  \item \texttt{ResendRequest} is parameterised by frame indices,
        and replay is read from a frame-addressable durable buffer.
\end{enumerate}
We bound this in Theorem~2 of the companion proof document.

\section{Compliance and Audit}
\label{sec:compliance}

\paragraph{MiFID~II RTS~25.}
RTS~25 mandates clock synchronisation and timestamp granularity
appropriate to the asset class; for HFT-eligible flow the threshold
is $1~\mu\mathrm{s}$ with a divergence from UTC bounded by
$100~\mu\mathrm{s}$~\cite{rts25}. ZAP-FIX frames carry a transport
timestamp populated from a hardware \texttt{PTP/IEEE-1588}-disciplined
clock, signed inside the frame's authenticated additional data. The
signed transport timestamp is in addition to the FIX
\texttt{TransactTime} (tag~60) carried in the body; the two are
independently auditable.

\paragraph{FINRA~OATS, ATS-N, CAT.}
The FIX session-layer artefacts that these regimes require
(\texttt{ClOrdID}, \texttt{OrigClOrdID}, event timestamps,
\texttt{ExecID})~\cite{finra-oats,finra-cat} are unchanged by
ZAP-FIX; the binding adds, rather than replaces, audit signal.
Crucially, the ZAP frame's authenticated transcript is itself an
admissible audit artefact: the absence of a frame in the durable
log implies the absence of the corresponding FIX message at the
session layer, with cryptographic strength.

\paragraph{Post-quantum confidentiality.}
For asset classes whose order intent retains economic value over
multi-year horizons (some primary issuance, certain block trades,
strategic accumulation programmes), harvest-now-decrypt-later is a
realistic threat. ZAP-FIX inherits the X-Wing hybrid KEM~\cite{xwing}
of the underlying ZAP transport; FIX-over-TLS does not, in any
deployed profile.

\section{Compatibility}
\label{sec:compat}

ZAP-FIX coexists with deployed FIX infrastructure through a
transparent translation gateway. The gateway terminates a TCP/TLS-FIX
session on its public face and a ZAP-FIX session on its private
face. Because the FIX session layer is unchanged, the gateway is a
stateful but stateless-with-respect-to-application-semantics relay:
it forwards Logon, Logout, Heartbeat, and SequenceReset transparently,
and translates between the SOH byte stream and the ZAP envelope.

This means an institutional client whose order-management system is
certified against a specific FIX dialect over TCP+TLS need not
recertify or modify its endpoint. Migration to ZAP-FIX endpoints can
proceed at the operator's pace.

\section{Related Work}
\label{sec:related}

FIX~FAST~\cite{fixfast} compresses FIX market-data feeds via implicit
field presence and stop-bit encodings; it is widely deployed for
multicast market data but uncommon for order entry. FIX~SBE
(Simple Binary Encoding)~\cite{fixsbe} replaces SOH text with a
fixed-layout binary encoding; it presupposes a transport binding
and is most often layered over TCP. ZAP-FIX BinaryEncoding mode is
philosophically closest to SBE but inherits transport-layer security
and post-quantum guarantees from ZAP.

FIX-over-WebSocket extensions exist for browser-side market data and
retail-adjacent order entry; they trade efficiency for firewall
traversal. They are not used in colocated institutional channels
and are not a competitor in the regime addressed here.

The HFT-latency literature gives detailed budgets and tail-latency
analyses against which our model can be
calibrated~\cite{hasbrouck2013,kissell2013,latencylottery}.
Colocated venue architectures from Refinitiv, Bloomberg, and SFTI
universally expose TCP+TLS-FIX as the institutional binding;
ZAP-FIX is intended as a forward-compatible upgrade path.

\section{Conclusion}

ZAP-FIX preserves the FIX session protocol bit-for-bit while
substituting a zero-copy, mutually authenticated, post-quantum
transport for the TCP+TLS byte stream. The construction is
transparent to FIX session-layer logic, recoverable through the
existing GapFill/ResendRequest mechanisms, and strictly improves on
the steady-state and resumption latency of TCP+TLS-FIX. The
companion proof document gives the formal latency model.

\bibliographystyle{plain}
\bibliography{refs}

\end{document}
